\clearpage
% Numérotation romaine pour annexe et bibliographie
\pagenumbering{roman}
\setcounter{page}{11}

\chapter*{Annexe A}
\addcontentsline{toc}{chapter}{Annexe A}
\thispagestyle{plain}
\markboth{}{}
% Numérotation des tableaux de l'annexe : A.1, A.2, …
\setcounter{table}{0}
\renewcommand{\thetable}{A.\arabic{table}}

Cette annexe rassemble les données géométriques, matérielles et numériques qui complètent le Chapitre~III et permettent de reproduire le calcul.

\section*{A.1 Géométrie et matériel de référence}

\subsection*{A.1.1 Ancrages et attaches}

Le tableau~A.1 donne les coordonnées des ancrages $\mathbf{A}_i$, exprimées dans le repère fixe $R_0$, et celles des attaches $\mathbf{b}_i$, exprimées dans le repère de la plateforme ($300 \times 240$~mm, orientation fixe $R = I$). Les câbles 1 à 4 sont les câbles hauts et les câbles 5 à 8 les câbles bas. Les câbles 2 à 4 et 6 à 8 sont numérotés par symétrie à partir des câbles 1 et 5.

\begin{table}[H]
\centering
\caption{Coordonnées des ancrages $\mathbf{A}_i$ ($R_0$) et des attaches $\mathbf{b}_i$ (plateforme), en mètres}
\label{tab:A1}
\footnotesize
\begin{tabular}{|c|c|c|c|c|c|c|}
\hline
\textbf{Câble} & $\mathbf{A}_{ix}$ & $\mathbf{A}_{iy}$ & $\mathbf{A}_{iz}$ & $\mathbf{b}_{ix}$ & $\mathbf{b}_{iy}$ & $\mathbf{b}_{iz}$ \\
\hline
1 & $-1{,}25$ & $+1{,}25$ & $+1{,}00$ & $-0{,}150$ & $+0{,}120$ & $+0{,}025$ \\
\hline
2 & $+1{,}25$ & $+1{,}25$ & $+1{,}00$ & $+0{,}150$ & $+0{,}120$ & $+0{,}025$ \\
\hline
3 & $-1{,}25$ & $-1{,}25$ & $+1{,}00$ & $-0{,}150$ & $-0{,}120$ & $+0{,}025$ \\
\hline
4 & $+1{,}25$ & $-1{,}25$ & $+1{,}00$ & $+0{,}150$ & $-0{,}120$ & $+0{,}025$ \\
\hline
5 & $-1{,}25$ & $+1{,}25$ & $0{,}00$ & $-0{,}150$ & $+0{,}120$ & $-0{,}025$ \\
\hline
6 & $+1{,}25$ & $+1{,}25$ & $0{,}00$ & $+0{,}150$ & $+0{,}120$ & $-0{,}025$ \\
\hline
7 & $-1{,}25$ & $-1{,}25$ & $0{,}00$ & $-0{,}150$ & $-0{,}120$ & $-0{,}025$ \\
\hline
8 & $+1{,}25$ & $-1{,}25$ & $0{,}00$ & $+0{,}150$ & $-0{,}120$ & $-0{,}025$ \\
\hline
\end{tabular}
\end{table}

Pour une pose $\mathbf{p}$, l'attache exprimée dans $R_0$ vaut $\mathbf{B}_i(\mathbf{p}) = \mathbf{p} + \mathbf{b}_i$ et la longueur du câble est $l_i = \|\mathbf{A}_i - \mathbf{B}_i(\mathbf{p})\|$ (équation~III.1). À la position centrale $\mathbf{p} = (0\,;\, 0\,;\, 0{,}45)$~m, le câble~1 donne $l_1 = 1{,}662$~m, valeur reprise au tableau~III.4.

\subsection*{A.1.2 Actionneurs et transmission}

Le tableau~A.2 précise les composants qui interviennent dans la borne $T_{\max} = 58$~N du Chapitre~III.

\begin{table}[H]
\centering
\caption{Spécifications matérielles}
\label{tab:A2}
\footnotesize
\begin{tabular}{|p{3.6cm}|p{3.4cm}|p{5.2cm}|}
\hline
\textbf{Composant} & \textbf{Paramètre} & \textbf{Valeur} \\
\hline
Moteur pas à pas NEMA~17 & Couple de maintien & $0{,}40$ à $0{,}47$~N$\cdot$m \\
\cline{2-3}
 & Angle de pas & $1{,}8^\circ$ (200 pas/tour) \\
\cline{2-3}
 & Courant par phase & $1{,}5$ à $1{,}7$~A \\
\hline
Tambour & Diamètre / rayon & $80$~mm / $40$~mm \\
\cline{2-3}
 & Enroulement par tour & $0{,}2513$~m \\
\hline
Câble Dyneema & Diamètre & $4$~mm \\
\cline{2-3}
 & Module d'Young & $80$ à $120$~GPa \\
\cline{2-3}
 & Charge de rupture & $> 500$~N \\
\hline
Poulie & Type / rayon & articulée (type WZ~31) / $20$~mm \\
\hline
\end{tabular}
\end{table}

\paragraph{Justification de $T_{\max}$.}
La tension qu'un actionneur peut exercer vaut $T = C / r$, où $C$ est le couple disponible au tambour et $r$ son rayon. Sur un tambour de $40$~mm et sans réduction, un couple de $0{,}47$~N$\cdot$m donne $T \approx 11{,}8$~N. La borne $T_{\max} = 58$~N retenue pour le calcul suppose donc une réduction d'environ $5{:}1$ entre le moteur et le tambour. Elle est utilisée comme borne supérieure de dimensionnement. Le calcul de référence n'en dépend pas : les tensions maximales obtenues ($4{,}395$~N) restent inférieures à la capacité du moteur en entraînement direct.

Le rapport entre la charge de rupture du câble et $T_{\max}$ dépasse $8$, ce qui laisse une large marge de sécurité. Le modèle de câble reste idéal (inextensible, sans masse) : l'élasticité, la flèche et le frottement de poulie ne sont pas pris en compte (voir~A.7).

\section*{A.2 Configuration du calcul et grille des points de vue}

\subsection*{A.2.1 Commande de référence}

Tous les résultats du Chapitre~III proviennent d'un seul calcul :
\begin{verbatim}
python scripts/run_pipeline.py --n 7 --zero-loss --pop-size 500 \
    --n-gen 5000 --output-dir data/outputs_7x7_pop500_gen5000 --verbose
\end{verbatim}

Les options de la commande priment sur la configuration par défaut : la grille est de $7\times7$ ($N = 49$). La graine du générateur aléatoire est $42$, ce qui rend l'optimisation NSGA-II déterministe. Les fichiers produits (CSV et JSON) sont conservés dans l'archive du calcul et se régénèrent avec cette commande.

\begin{table}[H]
\centering
\caption{Paramètres complémentaires du calcul}
\label{tab:A3}
\footnotesize
\begin{tabular}{|l|p{4.0cm}|p{6.5cm}|}
\hline
\textbf{Famille} & \textbf{Paramètre} & \textbf{Valeur} \\
\hline
Mission & Hauteur de prise de vue & $0{,}45$~m dans $R_0$ ($z_{\mathrm{vol}} = 0{,}30$~m, plants $0{,}35$~m, offset sol $0{,}20$~m) \\
\hline
Mission & Stabilisation par point & $2$~s ($49 \times 2 = 98$~s au total) \\
\hline
Vision & Champ de vue de la caméra & $62{,}2^\circ \times 48{,}8^\circ$ \\
\hline
Cinématique & $a_{\max}$ & $0{,}030$~m/s$^2$ \\
\hline
Cinématique & $T_{\mathrm{acc}} / T_{\mathrm{dec}}$ (profil LSPB) & $1{,}5$~s / $1{,}5$~s \\
\hline
NSGA-II & Opérateurs & croisement OX, mutation par inversion \\
\hline
\end{tabular}
\end{table}

\subsection*{A.2.2 Grille des points de vue}

Les $49$ points sont répartis uniformément dans le plan $z = 0{,}45$~m :
\begin{equation*}
x_k = -0{,}400 + k \times 0{,}1333~\text{m},
\qquad
y_m = -0{,}240 + m \times 0{,}0800~\text{m},
\quad k, m \in \{0, \ldots, 6\}.
\end{equation*}
Le pas est donc de $0{,}1333$~m selon $X$ et de $0{,}0800$~m selon $Y$. Le point~$0$ est $(-0{,}400\,;\, -0{,}240\,;\, 0{,}450)$. Le centre (id~$24$), point de départ et de retour de la mission, est $(0\,;\, 0\,;\, 0{,}450)$. Le point~$48$ est $(+0{,}400\,;\, +0{,}240\,;\, 0{,}450)$.

L'empreinte de la caméra à $0{,}45$~m est d'environ $0{,}3619 \times 0{,}2722$~m. Les recouvrements entre points voisins sont de $0{,}2286$~m en $X$ et $0{,}1922$~m en $Y$, ce qui assure une couverture complète. Les $49$ points sont tous retenus par le filtrage (aucun rejet, pour un seuil configuré de $15\,\%$).

\section*{A.3 Hypothèses de modélisation}

Cinq hypothèses ont été figées pour le calcul.

\begin{table}[H]
\centering
\caption{Hypothèses B1 à B5}
\label{tab:A4}
\footnotesize
\begin{tabular}{|c|p{3.8cm}|p{2.8cm}|p{4.5cm}|}
\hline
\textbf{Code} & \textbf{Hypothèse} & \textbf{Valeur} & \textbf{Statut} \\
\hline
B1 & Masse opérationnelle & $0{,}383988$~kg & Estimation calibrée ; à confirmer par pesée \\
\hline
B2 & Hauteur des plants & $0{,}35$~m & Estimation pour un haricot nain ; à mesurer \\
\hline
B3 & Champ de vue de la caméra & $62{,}2^\circ \times 48{,}8^\circ$ & Valeur constructeur \\
\hline
B4 & Grille et effectif & $7\times7$ ; $N = 49$ & Validée par le WFW sur les $49$ points \\
\hline
B5 & Offset sol / $R_0$ & $0{,}20$~m & Confirmé par le montage \\
\hline
\end{tabular}
\end{table}

\section*{A.4 Statistiques de la matrice de coûts}

La matrice contient $49 \times 49 = 2401$ entrées. La diagonale est nulle, toutes les paires sont réalisables et les valeurs sont symétriques dans les tolérances du validateur.

\begin{table}[H]
\centering
\caption{Statistiques (toutes paires ordonnées, diagonale incluse)}
\label{tab:A5}
\begin{tabular}{|l|c|c|c|}
\hline
\textbf{Statistique} & \textbf{Distance (m)} & \textbf{Temps (s)} & \textbf{Énergie (J)} \\
\hline
Minimum & $0{,}000$ & $0{,}000$ & $0{,}000$ \\
\hline
Médiane & $0{,}400$ & $10{,}389$ & $14{,}612$ \\
\hline
Moyenne & $0{,}391$ & $10{,}163$ & $15{,}431$ \\
\hline
Maximum & $0{,}933$ & $22{,}232$ & $49{,}240$ \\
\hline
\end{tabular}
\end{table}

Ces valeurs concernent des déplacements élémentaires entre deux points. Elles ne correspondent pas aux cumuls d'une trajectoire complète, qui dépendent de l'ordre de visite.

\section*{A.5 Validation logicielle}

Six contrôles sur sept sont réussis, un donne un avertissement et aucun n'échoue (tableau~A.6).

\begin{table}[H]
\centering
\caption{Résultats de validation}
\label{tab:A6}
\footnotesize
\begin{tabular}{|p{3.6cm}|c|p{7.2cm}|}
\hline
\textbf{Contrôle} & \textbf{Statut} & \textbf{Évidence} \\
\hline
Cohérence de la configuration & PASS & $N = 49$ ; $T \in [1\,;\, 58]$~N ; orientation fixe \\
\hline
Points filtrés & PASS & $49$ points ; aucun infaisable \\
\hline
Matrice de coûts & PASS & $2401$ entrées ; symétrie ; diagonale nulle \\
\hline
Solutions de Pareto & WARN & Une solution non dominée unique ; seuil recommandé $\geq 10$ \\
\hline
Trajectoires & PASS & $3$ types ; centre aux extrémités ; totaux cohérents \\
\hline
Profils cinématiques & PASS & $5361$ échantillons ; $\|\mathbf{v}\|_{\max} = 0{,}045$~m/s \\
\hline
Cohérence croisée & PASS & Profils, matrice et résumé concordants \\
\hline
\end{tabular}
\end{table}

\paragraph{Interprétation de l'avertissement.}
Le fichier \texttt{solutions\_pareto.csv} contient deux ordres de visite distincts dont les objectifs sont rigoureusement identiques : ($177{,}295$~s ; $246{,}436$~J). Le validateur les compte comme une seule solution non dominée, d'où $n_{\mathrm{solutions}} = 1$. Il s'agit d'un doublon de représentation du même point objectif, et non d'une incohérence. L'avertissement signale seulement que le front est trop réduit pour étudier un compromis entre temps et énergie.

\section*{A.6 Choix de modélisation}

\paragraph{Convention de signe.}
La direction unitaire est $\mathbf{U}_i = (\mathbf{A}_i - \mathbf{B}_i) / \|\mathbf{A}_i - \mathbf{B}_i\|$, orientée de l'attache vers l'ancrage, c'est-à-dire dans le sens physique de la traction. La matrice de structure $W$ assemble les vecteurs $[\mathbf{U}_i\,;\, \mathbf{B}_i \times \mathbf{U}_i]$. Cette convention évite d'obtenir des tensions positives qui correspondraient en réalité à une compression de la plateforme.

\paragraph{Résolution des tensions en deux chemins.}
À chaque sous-pas, une solution par pseudo-inverse est tentée en premier. Si elle sort de $[T_{\min}\,;\, T_{\max}]$, un solveur borné SLSQP, initialisé par la solution précédente, prend le relais. Le coût de calcul moyen est ainsi réduit sans perdre la robustesse dans les cas limites.

\paragraph{Énergie au point milieu.}
L'énergie de chaque sous-segment est évaluée en son point milieu, ce qui rend $E(i,j)$ et $E(j,i)$ symétriques par construction.

\section*{A.7 Limites et protocole expérimental}

\begin{table}[H]
\centering
\caption{Limites et actions associées}
\label{tab:A7}
\footnotesize
\begin{tabular}{|p{3.8cm}|p{4.2cm}|p{4.5cm}|}
\hline
\textbf{Limite} & \textbf{Impact} & \textbf{Action} \\
\hline
Masse estimée (B1) & Influence le poids et les tensions & Peser la plateforme ; relancer le calcul si l'écart dépasse $5\,\%$ \\
\hline
$T_{\max} = 58$~N & Dépend du couple, de la réduction et du rayon du tambour & Vérifier la chaîne de transmission complète \\
\hline
Absence de vent et de perturbations & Efforts externes non modélisés & Ajouter des scénarios de charge et de replanification \\
\hline
Solution non dominée unique & Pas d'étude statistique du compromis & Répéter avec plusieurs graines \\
\hline
Câble idéal & Élasticité, flèche et frottement négligés & Valider par mesure des tensions \\
\hline
Stabilisation constante ($2$~s) & Indépendante du mouvement et des vibrations & Mesurer le temps réel sur prototype \\
\hline
\end{tabular}
\end{table}

